% !TEX root = thesis.tex
\chapter{Experiments}\label{sec:experiments}

In this chapter, we present the experimental parameters and performance results based on the framework and downstream applications described in the previous chapter. We compare the raw CSI baseline with four distinct representations:

\begin{itemize}
    \item \textbf{CLS Embeddings:} The global representation vector extracted from the LWM transformer encoder.
    \item \textbf{Channel Embeddings:} The token-level spatial-frequency representation.
    \item \textbf{AE - 1/16:} A compressed vector $\mathbf{z}$ generated by an Autoencoder with a 1:16 compression ratio.
    \item \textbf{AE - 1/32:} A compressed vector $\mathbf{z}$ generated by an Autoencoder with a 1:32 compression ratio.
\end{itemize}

Furthermore, for the CLS representation, we compare the performance of the frozen backbone against partial fine-tuning (the last three layers) and full fine-tuning.

Since one of the primary goals of this work is to demonstrate that the compressed vector $\mathbf{z}$ can be used efficiently for downstream tasks, we only utilized the encoder portion of the autoencoder during evaluation. All calculated inference times, parameter counts, and floating-point operations (FLOPs) strictly reflect this architectural choice.

For all downstream tasks, synthetic data was generated using real-life scenarios from the DeepMIMO dataset \cite{Alkhateeb2019}. The Denver scenario (\texttt{city\_18\_denver}) features 3 available BSs and 6,970 User Equipments (UEs) uniformly distributed across the user grid. These maps, illustrated in \ref{fig:scenario_maps}, show the urban topology and the coverage map for a selected BS. Because this specific scenario was not seen by the LWM or the Autoencoders during their initial pretraining, it serves as an excellent benchmark to showcase the models' generalization capabilities in a standard urban environment.

\begin{figure}[htbp]
    \centering
    \begin{subfigure}{0.4\textwidth}
        \centering
        \includegraphics[width=\textwidth]{experiments/city_18_denver_map.png}
        \caption{Denver (\texttt{city\_18\_denver})}
        \label{fig:map_denver}
    \end{subfigure}
    \hfill
    \begin{subfigure}{0.55\textwidth}
        \centering
        \includegraphics[width=\textwidth]{experiments/city_18_denver_coverage_map.pdf}
        \caption{Coverage map for different selected BSs.}
        \label{fig:city_18_denver_coverage_map}
    \end{subfigure}
    \caption{DeepMIMO ray-tracing simulation environment used for dataset generation.}
    \label{fig:scenario_maps}
\end{figure}


The dataset generation uses a standardized configuration across all scenarios to ensure a consistent comparison. At the BS, a ULA is employed, consisting of 32 antenna elements with a predefined rotation angle of $-135^\circ$ in the horizontal plane. On the user side, every UE is equipped with a single antenna. Furthermore, the ray-tracing engine computes 20 distinct multipath components for every channel link.


In the frequency domain, the system employs Orthogonal Frequency-Division Multiplexing (OFDM) with 32 active subcarriers and a subcarrier spacing of 30 kHz. Additionally, in the Denver simulation, we activate all three BSs and filter out any samples that lacked a valid path between the UE and the BS. However, for the specific downstream task of power allocation, UEs were grouped strictly within the coverage area of their respective serving BS.

\begin{table}[htbp]
    \centering
    \caption{DeepMIMO Dataset Generation Parameters}
    \label{tab:deepmimo_generation_params}
    \begin{tabular}{ll}
        \hline
        \textbf{Parameter} & \textbf{Value} \\
        \hline
        BS Antenna Array Shape ($N$) & $32 \times 1$ (Uniform Linear Array) \\
        BS Antenna Rotation ($x, y, z$) & $(0^\circ, 0^\circ, -135^\circ)$ \\
        UE Antenna Array Shape & $1 \times 1$ (Single Antenna) \\
        Number of Multipath Components & 20 \\
        Number of Subcarriers ($M$) & 32 \\
        Subcarrier Spacing (SCS) & 30 kHz \\
        Total Bandwidth & 960 kHz \\
        Active Base Station (BS) Index & 1, 2, 3 (Denver) \\
        \hline
    \end{tabular}
\end{table}

Once the channel matrices $\mathbf{H}$ between each user and the selected BSs are generated, the complex coefficients are multiplied by a scaling factor of $\alpha = 10^{6}$. This uniform scaling strategy was adopted to increase training stability and ensure that the input magnitudes align with the distribution expected by the pretrained LWM model. This adjustment minimizes the risk of vanishing gradients and allows for more stable downstream fine-tuning. For all experiments, the dataset was strictly partitioned into 60\% for training, 20\% for validation, and 20\% for testing.

For every experiment, we systematically evaluate the learned representations across three dimensions: data efficiency, noise robustness, and computational complexity. Furthermore, we perform an analysis of the latent space through t-Distributed Stochastic Neighbor Embedding (t-SNE) visualizations \cite{vandermaaten2008visualizing}. The goal of this qualitative analysis is to demonstrate that the pretrained models are capable of inherently capturing channel conditions and spatial features within their latent manifolds.

Additionally, we measured the number of parameters and FLOPs utilized by the LWM and the Autoencoders. As summarized in \ref{tab:encoder_complexity}, despite requiring fewer overall FLOPs, the LWM exhibits a higher memory footprint than the Autoencoders. This discrepancy is explained by the fundamental architectural differences between the sequential, highly parallelizable operations of CNNs and the memory-bandwidth-bound self-attention mechanisms of Transformers.

\begin{table}[htbp]
    \centering
    \caption{Computational Complexity of the Encoders}
    \label{tab:encoder_complexity}
    \begin{tabular}{lll}
        \hline
        \textbf{Model} & \textbf{Parameters (M)} & \textbf{MFLOPs} \\
        \hline
        LWM & 2.47 & 6.65 \\
        AE - 1/16 & 1.16 & 59.02 \\
        AE - 1/32 & 0.64 & 58.49 \\
        \hline
    \end{tabular}
\end{table}

Finally, all training loops and benchmark tests were conducted using an NVIDIA L4 GPU provisioned through Google Colab.

\section{Line-of-Sight (LoS) Classification Results}\label{sec:experiments_los_nlos}

The first set of experiments focuses on the binary classification of LoS and NLoS conditions.

For this downstream task, we utilized the downstream model (or Task Head) described in Section \ref{sec:task_los_nlos_model} with a hidden dimension equal to 256. As stated in the previous chapter, the input is the flattened output of the Autoencoder, the flattened LWM embeddings, or the raw channel matrix. Since this is a binary classification task, the final output size is always 2. Because of the varying input dimensions, the number of parameters and FLOPs of the downstream model depends heavily on the chosen representation.

\begin{table}[htbp]
    \centering
    \caption{Computational Complexity of the Downstream Model}
    \label{tab:los_nlos_model_complexity}
    \begin{tabular}{lll}
        \hline
        \textbf{Input Type} & \textbf{Head Parameters (K)} & \textbf{Head KFLOPs} \\
        \hline
        Raw Channel & 525 & 525 \\
        LWM - CLS & 33 & 33 \\
        LWM - Channel Embeddings & 2097 & 2097 \\
        AE - 1/16 & 34 & 34 \\
        AE - 1/32 & 17 & 17 \\
        \hline
    \end{tabular}
\end{table}

The \ref{tab:los_nlos_model_complexity} presents the full number of parameters and flops for the downstream model used in each representation.

\subsection{Data Efficiency}\label{sec:experiments_los_nlos_data_eff}

To evaluate the data efficiency of different representations, we trained the downstream MLP using a wide range of training volumes. Specifically, we sampled from 1\% of the training pool, which corresponds to 23 samples in the Denver dataset, up to 100\% of the available training data (2,355 samples).

\begin{figure}[htbp]
    \centering
    \includegraphics[width=0.5\textwidth]{experiments/los_nlos/data_eff.pdf}
    \caption{Data efficiency results for LoS/NLoS classification in the Denver scenario.}
    \label{fig:data_eff_los_nlos_denver}
\end{figure}

The \ref{fig:data_eff_los_nlos_denver} shows that the learned representations perform well in this scenario, with the Channel Embeddings demonstrating excellent performance in the few-shot regime. These representations saturate significantly earlier than the Raw CSI baseline. However, given enough data, all representations eventually converge to similar performance levels.

\begin{figure}[htbp]
    \centering
    \includegraphics[width=0.5\textwidth]{experiments/los_nlos/data_eff_fine_tune.pdf}
    \caption{Impact of fine-tuning on LoS/NLoS Classification performance.}
    \label{fig:data_eff_fine_tune_los_nlos_denver}
\end{figure}

These results demonstrate a clear advantage in low-data environments, where the representations learned by the LWM and the Autoencoders significantly outperform the baseline. Nevertheless, while the global CLS representation can struggle in certain extreme low-data scenarios, its performance can be vastly improved through fine-tuning as seen in \ref{fig:data_eff_fine_tune_los_nlos_denver}.


\subsection{Noise Robustness}\label{sec:experiments_los_nlos_noise}

To evaluate noise robustness, we trained the downstream models using 100\% of the training pool. During the testing phase, we dynamically applied AWGN to the channel samples to match desired Signal-to-Noise Ratio (SNR) targets. The result is shown as \ref{fig:noise_robustness_los_nlos_denver}.

\begin{figure}[htbp]
    \centering
    \includegraphics[width=0.5\textwidth]{experiments/los_nlos/noise_robustness.pdf}
    \caption{Effect of noise on the F1-Score using the full training dataset.}
    \label{fig:noise_robustness_los_nlos_denver}
\end{figure}

In extremely low SNR regimes, all representations exhibited significant performance degradation. However, the compressed vectors generated by the Autoencoders demonstrated superior robustness compared to the LWM embeddings. 

\begin{figure}[htbp]
    \centering
    \includegraphics[width=0.5\textwidth]{experiments/los_nlos/noise_robustness_fine_tune.pdf}
    \caption{Effect of fine-tuning on noise robustness measured by F1-Score.}
    \label{fig:noise_robustness_fine_tune_denver}
\end{figure}

The \ref{fig:noise_robustness_fine_tune_denver} shows that while full fine-tuning can marginally improve the noise robustness of the CLS representation, partial fine-tuning failed to do so in our experiments. This counterintuitive behavior likely occurs because the model overfits to the clean representations.

\subsection{Computational Complexity}\label{sec:experiments_los_nlos_complexity}

Latency was calculated by averaging the execution time of 500 inference cycles per sample. As shown in Table \ref{tab:latency_analysis}, the head latency remained relatively constant across the different representations. However, it is clear that the LWM incurs a significant computational cost in the encoder phase.

\begin{table}[htbp]
    \centering
    \caption{Inference Time}
    \label{tab:latency_analysis}
    \begin{tabular}{lccccc}
        \hline
        \textbf{Model Name} & \textbf{Encoder (ms)} & \textbf{Head (ms)} & \textbf{Total Time (ms)} \\
        \hline
        Raw & - & 0.17 & 0.17 \\
        LWM - CLS & 10.78 & 0.20 & 10.98 \\
        LWM - Channel Embeddings & 10.89 & 0.20 & 11.09 \\
        AE - 1/16 & 0.73 & 0.18 & 0.91 \\
        AE - 1/32 & 0.74 & 0.18 & 0.92 \\
        \hline
    \end{tabular}
\end{table}

Additionally, the training time was measured to understand the overhead introduced during model optimization. The analysis in \ref{tab:training_analysis} demonstrates that while the Autoencoders add negligible training overhead compared to the baseline, the LWM adds a substantial computational burden. This cost is further exacerbated when fine-tuning is applied.

\begin{table}[htbp]
    \centering
    \caption{Training Time Analysis}
    \label{tab:training_analysis}
    \begin{tabular}{lc}
        \hline
        \textbf{Model Name} & \textbf{Training Time (s)} \\
        \hline
        Raw & 1.90 \\
        LWM - CLS & 15.17 \\
        LWM - CLS (Partial Fine-Tuning) & 24.56 \\
        LWM - CLS (Full Fine-Tuning) & 54.42 \\
        LWM - Channel Embeddings & 15.51 \\
        AE - 1/16 & 2.83 \\
        AE - 1/32 & 2.85 \\
        \hline
    \end{tabular}
\end{table}

\subsection{Latent Space Analysis}\label{sec:experiments_los_nlos_latent}

To understand the underlying mechanisms driving the superior data efficiency and noise robustness of the learned representations, we conducted a qualitative analysis using t-SNE. By projecting the high-dimensional embeddings into a 2D space, we can visualize how the models intrinsically cluster the channel conditions.

\begin{figure}[htbp]
    \centering
    \includegraphics[width=\linewidth]{experiments/los_nlos/tsne_comparison.pdf}
    \caption{Latent space visualizations for the evaluated representations. The figure is complemented by the scenario layout.}
    \label{fig:tsne_comparison_data_los_nlos_denver}
\end{figure}

As seen in \ref{fig:tsne_comparison_data_los_nlos_denver}, the t-SNE projections provide compelling visual evidence to support the preceding quantitative analyses:

\begin{itemize}
    \item The \textbf{Raw CSI} exhibits highly overlapping regions, which fundamentally complicates channel state classification for a simple downstream MLP.
    \item The \textbf{Autoencoders} demonstrate a greater degree of spatial separation, resulting in fewer overlapping boundaries between the classes.
    \item Even without fine-tuning, the \textbf{frozen LWM} exhibits an incredible degree of natural clustering, which explains the high data efficiency of this representation in the few-shot regime.
    \item After \textbf{full fine-tuning}, the embedding space becomes remarkably segregated. The LoS and NLoS conditions are almost perfectly linearly separated. For this reason, the downstream classification head is able to achieve a near-perfect F1-Score with only a trivial volume of training data.
\end{itemize}

\section{Beam Selection Results}\label{sec:experiments_beam}

The second set of experiments focuses on multi-class classification and best beam prediction. In this task, we added a complexity component related to the number of available beams. The downstream architecture was described in \ref{sec:task_beam_model}. For this experiment, we used a hidden dimension of 512. The input size, as in the previous downstream task, depends on the input type, while the output size is strictly equal to the number of beams.


As in the previous task, we measured the number of parameters and FLOPs to assess computational overhead and compiled the results at \ref{tab:beam_selection_model_complexity}.

\begin{table}[htbp]
    \centering
    \caption{Computational Complexity of the Downstream Model}
    \label{tab:beam_selection_model_complexity}
    \begin{tabular}{lll}
        \hline
        \textbf{Input Type} & \textbf{Head Parameters (K)} & \textbf{Head KFLOPs} \\
        \hline
        Raw Channel & 1215 & 1215 \\
        LWM - CLS & 232  & 232 \\
        LWM - Channel Embeddings & 4361 & 4361  \\
        AE - 1/16 & 232 & 232 \\
        AE - 1/32 & 200 & 200 \\
        \hline
    \end{tabular}
\end{table}

\subsection{Data Efficiency}\label{sec:experiments_beam_data_eff}

Similar to the previous task, we analyse how the data volume impacts overall performance.

\begin{figure}[h]
    \centering
    \includegraphics[width=\linewidth]{experiments/beam_selection/data_eff_beam.pdf}
    \caption{Heatmap visualizations of F1-Score with different data volumes and number of beams.}
    \label{fig:data_eff_beam_selection_denver}
\end{figure}

In this particular task, the data at \ref{fig:data_eff_beam_selection_denver} does not show a distinct advantage when using the learned representations. The results obtained with the compressed Autoencoder vector and the channel embeddings are very similar to the baseline. However, the global CLS embedding exhibits inferior performance in both the few-shot regime and high data volumes. All models begin to struggle as the task complexity increases. The compressed vector $\mathbf{z}$ performs similarly to the channel embeddings and the raw CSI despite its low-dimensional nature. For this reason, utilizing $\mathbf{z}$ is highly advantageous for multiple downstream tasks, as it reduces computational overhead and conserves transmission bandwidth.

\begin{figure}[h]
    \centering
    \includegraphics[width=\linewidth]{experiments/beam_selection/data_eff_fine_tune.pdf}
    \caption{Heatmap visualizations of the effect of fine-tuning on F1-Score with different data volumes and number of beams.}
    \label{fig:data_eff_fine_tune_beam_selection_denver}
\end{figure}

Even with fine-tuning (\ref{fig:data_eff_fine_tune_beam_selection_denver}), the CSI baseline outperforms the CLS embeddings. In very low data volumes (10\%), there is a slight benefit compared to the baseline; however, as the data volume grows, the CSI performance overcomes the global embeddings. This suggests the global CLS token lacks the micro-spatial relationships strictly needed for precise beam selection.

\subsection{Noise Robustness}\label{sec:experiments_beam_noise}

Under noisy conditions, the Raw CSI remains the most robust representation. The CLS representation's performance collapses even in the least complex scenarios. The compressed vector $\mathbf{z}$ demonstrates the strongest robustness among the learned representations; however, in extremely noisy conditions, it still lags behind the baseline. The channel embeddings show good results under low complexity and high SNR, but when one of these variables changes, the performance collapses rapidly.

\begin{figure}[h]
    \centering
    \includegraphics[width=\linewidth]{experiments/beam_selection/noise_robustness.pdf}
    \caption{Heatmap visualizations of the effect of noise on F1-Score with different number of beams.}
    \label{fig:noise_robustness_beam_selection_denver}
\end{figure}

Fine-tuning brought limited gains to noise resilience. Under perfect channel conditions, the performance moved closer to the baseline. However, partial fine-tuning was entirely incapable of improving noise robustness. Additionally, full fine-tuning, as observed in the previous task, yielded greater performance gains. This further indicates that the global representation learned by the LWM overfits to clean channels, making its performance difficult to generalize in noisy environments.

\begin{figure}[h]
    \centering
    \includegraphics[width=\linewidth]{experiments/beam_selection/noise_robustness_fine_tune.pdf}
    \caption{Heatmap visualizations of the effect of fine-tuning on F1-Score under noise with different number of beams.}
    \label{fig:noise_robustness_fine_tune_beam_selection_denver}
\end{figure}

\subsection{Computational Complexity}\label{sec:experiments_beam_complexity}

Using the same methodology as before, the latency and training time for 100\% of the samples were evaluated. The inference time of the encoder remained unchanged, but the inference time of the downstream head increased. This occurs due to the higher number of dense layers required for this specific multi-class downstream task.

\begin{table}[htbp]
    \centering
    \caption{Inference Time}
    \label{tab:latency_analysis_beam_selection}
    \begin{tabular}{lccccc}
        \hline
        \textbf{Model Name} & \textbf{Encoder (ms)} & \textbf{Head (ms)} & \textbf{Total Time (ms)} \\
        \hline
        Raw & - & 0.52 & 0.52 \\
        LWM - CLS & 10.76 & 0.58 & 11.34 \\
        LWM - Channel Embeddings & 10.91 & 0.59 & 11.50 \\
        AE - 1/16 & 0.72 & 0.53 & 1.25 \\
        AE - 1/32 & 0.70 & 0.52 & 1.22 \\
        \hline
    \end{tabular}
\end{table}

As in the previous task, the Autoencoder added minimal latency overhead. On the other hand, the use of the LWM, especially during fine-tuning, added significant computational costs to the training pipeline, which is verified in \ref{tab:training_analysis_beam_selection}.

\begin{table}[htbp]
    \centering
    \caption{Training Time Analysis (16 Beams)}
    \label{tab:training_analysis_beam_selection}
    \begin{tabular}{lc}
        \hline
        \textbf{Model Name} & \textbf{Training Time (s)} \\
        \hline
        Raw & 5.85 \\
        LWM - CLS & 30.03 \\
        LWM - CLS (Partial Fine-Tuning) & 47.00 \\
        LWM - CLS (Full Fine-Tuning) & 100.06 \\
        LWM - Channel Embeddings & 31.11 \\
        AE - 1/16 & 7.75 \\
        AE - 1/32 & 7.83 \\
        \hline
    \end{tabular}
\end{table}

\subsection{Latent Space Analysis}\label{sec:experiments_beam_latent}

To further understand the underlying representations, we performed a qualitative analysis through the visualization of the latent space using t-SNE. The \ref{fig:tsne_beam_selection} shows the behavior of the models for the selection problem using 16 beams. The data points are colored by their optimal beam index.


\begin{figure}[htbp]
    \centering
    \includegraphics[width=\linewidth]{experiments/beam_selection/tsne_comparison.pdf}
    \caption{t-SNE visualizations of the latent space for the beam selection task. The points are colored by their optimal beam index, illustrating the degree of spatial-feature entanglement across different representations.}
    \label{fig:tsne_beam_selection}
\end{figure}

The visual evidence strongly corroborates the performance trends observed in the preceding analysis:

\begin{itemize}
    \item The \textbf{CLS embeddings} have a highly entangled latent space. The different beam classes overlap heavily. This is expected because the global token was designed to capture macroscopic channel properties; it lacks the fine-grained magnitude and phase information necessary for this precise directional task.
    \item \textbf{Fine-tuning} completely transforms the latent space and creates a clear separation between different classes, which explains the performance gains over the frozen model.
    \item \textbf{Channel embeddings} produce an entangled latent space. However, their downstream classification performance is satisfactory, which indicates that the high-dimensional embeddings successfully extract the necessary spatial information even if it is not linearly separable in 2D space.
    \item The \textbf{Autoencoders} show moderate clustering, resulting in good performance compared to the LWM embeddings. This indicates that the convolutional network successfully compresses the channel state information while retaining the crucial spatial structure.
    \item While the \textbf{Raw Channel} does not exhibit natural clustering in the t-SNE projection, the data natively retains all the necessary physical information for beam selection, explaining its superior quantitative results.
\end{itemize}

\section{Power Allocation Results}\label{sec:experiments_power}

This task aims to evaluate if the learned representations are useful for tasks where it is necessary to manage the inter-user interference. For all the learned representations, we used the downstream model described on \ref{sec:task_power_model} with 256 hidden channels in the entry block and the two residual blocks. The input size depends on the representation and the output size is equal to the number of users, which is the complexity variable for this task.

For this task, we did not compare the fine-tuning because in every case it did not bring significant difference in the experiments.

In the densest simulation, 16 users, we measured the number of parameters and FLOPs used by the downstream model and generated the \ref{tab:experiments_power_model_complexity}.

\begin{table}[htbp]
    \centering
    \caption{Computational Complexity of the Downstream Model}
    \label{tab:experiments_power_model_complexity}
    \begin{tabular}{lll}
        \hline
        \textbf{Input Type} & \textbf{Head Parameters (K)} & \textbf{Head KFLOPs} \\
        \hline
        Raw Channel & 196 & 3147 \\
        LWM - CLS & 12  & 199 \\
        LWM - Channel Embeddings & 786 &  12,585 \\
        AE - 1/16 & 12 & 199 \\
        AE - 1/32 & 6 & 100 \\
        \hline
    \end{tabular}
\end{table}


\subsection{Data Efficiency}\label{sec:experiments_power_data_eff}

\begin{figure}[htbp]
    \centering
    \includegraphics[width=\textwidth]{experiments/power_allocation/data_eff.pdf}
    \caption{Data efficiency results for Power Allocation in the Denver scenario.}
    \label{fig:power_allocation_data_eff}
\end{figure}

During the tests, it was clear that the autoencoder representation had a better performance. Even in low-data volumes, it showed near optimal performance. At 2 and 4 users, the raw channels baseline did not match the autoencoder even with 100\% of the dataset. Both LWM representations had unsatisfactory results. They failed to improve as the volume of data increases and had an unstable training.

As the number of users increases, the theoretical upper bound for the sum-rate improves. This phenomenon happens because the base station has a greater user diversity to choose from. When the pool increases, the system has a higher probability of finding a set of users that has good channel conditions and low interference. For those reasons, the gap between different representations also vanishes. This condition will happen until the number of users becomes the same as the number of antennas in the BS.

\subsection{Noise robustness}\label{sec:experiments_power_noise}

\begin{figure}[htbp]
    \centering
    \includegraphics[width=\textwidth]{experiments/power_allocation/noise_robustness.pdf}
    \caption{Noise robustness results for Power Allocation in the Denver scenario.}
    \label{fig:power_allocation_noise}
\end{figure}

This experiment evaluates the noise robustness of the downstream model by scaling the SNR from -5 dB to 25 dB in four different complexity scenarios. It is worth noting that, as the number of users and the quality of the channel increases, the gap between the upper and lower bound also increases. This explains the importance of a strategy other than the naive equal-power allocation.

The compressed vector generated by the autoencoder showed good results in every scenario, especially with 2 and 4 users. Also, the CLS embeddings had a surprisingly competitive performance in the scenario with 16 users, where it almost matched the upper bound.

It is important to notice that the performance of all representations degrades as the channel conditions worsen. However, this fact is masked because the upper and lower bounds also degrade.

\subsection{Computational Complexity}\label{sec:experiments_power_complexity}

Again, we measured the computational complexity, latency and training time for 100\% of  the training samples. The inference time of the encoder is unchanged and the inference time of the head does not vary a lot across representation. The dominant factor is the time necessary to generate the embeddings

\begin{table}[htbp]
    \centering
    \caption{Inference Time}
    \label{tab:latency_analysis_power}
    \begin{tabular}{lccccc}
        \hline
        \textbf{Model Name} & \textbf{Encoder (ms)} & \textbf{Head (ms)} & \textbf{Total Time (ms)} \\
        \hline
        Raw & - & 1.30 & 1.30 \\
        LWM - CLS & 11.20 & 0.58 & 11.78 \\
        LWM - Channel Embeddings & 11.31 & 0.59 & 11.90 \\
        AE - 1/16 & 0.75 & 0.53 & 1.28 \\
        AE - 1/32 & 0.74 & 0.52 & 1.26 \\
        \hline
    \end{tabular}
\end{table}

The training time was measured for 16 users and compiled at \ref{tab:training_analysis_power}.

\begin{table}[htbp]
    \centering
    \caption{Training Time Analysis (16 Users)}
    \label{tab:training_analysis_power}
    \begin{tabular}{lc}
        \hline
        \textbf{Model Name} & \textbf{Training Time (s)} \\
        \hline
        Raw & 42.04 \\
        LWM - CLS & 431.91 \\
        LWM - Channel Embeddings & 436.01 \\
        AE - 1/16 & 90.33 \\
        AE - 1/32 & 90.00 \\
        \hline
    \end{tabular}
\end{table}

\subsection{Latent Space Analysis}\label{sec:experiments_power_latent}

\begin{figure}[htbp]
    \centering
    \includegraphics[width=\linewidth]{experiments/power_allocation/tsne_comparison.pdf}
    \caption{t-SNE visualizations of the latent space for the power allocation task with 4 users. The points are colored by the ordinal power ranking of the user. Users that receive the most power are assigned the same color.}
    \label{fig:tsne_power_allocation}
\end{figure}

The analysis at \ref{fig:tsne_power_allocation} shows that it is difficult to organize the power budget only through spatial segregation in the latent space. All the representations have a highly entangled space. The raw channel baseline reveals complete absence of organization.

The autoencoder was able to organize more the latent space. However, it does not have perfect boundaries. Finally, both LWM embeddings showed improvements compared to the other representations, but it was not enough to create a good topology.

This demonstrates that the power allocation is a more complex task. Since the output depends on the interference between users, the representation is not enough to predict the allocated power. For this reason, the task demanded a more powerful downstream model.